%! Author = trdng
%! Date = 3/19/2026

% Preamble
\documentclass[11pt]{article}

% Packages
\usepackage{amsmath}
\usepackage{braket}
\usepackage{indentfirst}

\usepackage{graphicx}
\usepackage{listings}
\usepackage{xcolor}

% Listings style
\lstset{
    basicstyle=\ttfamily\small,
    keywordstyle=\color{blue}\bfseries,
    commentstyle=\color{green!50!black}\itshape,
    stringstyle=\color{red},
    showstringspaces=false,
    breaklines=true,
    frame=single
}

% Document
\begin{document}

    \section*{Wave Function}

    \indent The state of a particle (qubit) is described by a wave function.
    For a single qubit, the wave function is written as
    \[
        \ket{\psi} = \alpha\ket{0} + \beta\ket{1}
    \]

    For two-qubit system, the wave function is written as

    \[
        \ket{\psi} = \alpha\ket{00} + \beta\ket{01} + \gamma\ket{10} + \sigma\ket{11}
    \]

    For $n$-qubit system, the wave function is written as

    \[
        \ket{\psi} = \sum_{i=0}^{2^n - 1}\alpha_i\ket{i}
    \]

    For any qubit, the state can be in superposition described via the wave function.
    Superposition means the qubit can be both 0 and 1 at the same time until measured.

    When you measure a qubit
    \begin{itemize}
        \item The superposition disappears.
        \item The state becomes either 0 or 1.
        \item The result is probabilistic.
    \end{itemize}

    The transition "before-after" measured is called wave function collapses, and it converts quantum information into a classical result.
    Quantum circuits often do, for example

    \[
        \ket{0} \xrightarrow{} operations \xrightarrow{} measure
    \]

    The superposition persists throughout the circuit.
    Only the final measurement collapses it.

    \newpage

    \begin{lstlisting}[language=Python, caption={Measurement procedure in single qubit system},label={lst:lstlisting22}, captionpos=b]
        import numpy as np

        psi = np.array([1/np.sqrt(2), 1/np.sqrt(2)])
        probabilities = np.abs(psi) ** 2
        result = np.random.choice(np.arange(len(psi)), p=probabilities)

        collapsed = np.zeros_like(psi)
        collapsed[result] = 1

        print(f'Measured: {result}')
        print(f'Collapsed State: {collapsed}')
    \end{lstlisting}


\end{document}